\chapter{Regularity of eigenfunctions}
\label{chap:regularity}

The odd--odd and mixed arguments need regularity well beyond
$\omega\in L^2$: weighted bounds at the saddles, $H^1$ regularity across the
separatrix and $H^2$ regularity of quotients $\psi/\Psi$ or $\psi/\cos a$.
The transport gains must apply to jet-subtracted pairs, not only to
eigenfunctions satisfying an elliptic relation. The final section records
the separate elliptic-center argument; neither sector exclusion needs it
once the mixed boundary trace is handled locally.
In this chapter $d$ is distance to the saddle set and $r$ distance to a fixed saddle.

\section{Weighted elliptic estimates at a saddle}

\begin{lemma}[Cylinder inverse]
\label{lem:cylinder-inverse}
\uses{def:sobolev-local}
Let $u\in H^2$ on a disk about a saddle, with a prescribed reflection parity
in two orthonormal separatrix coordinates, and $\Delta u=F$. Let $a>1$
avoid the nonnegative angular frequencies allowed by the parity. If
$r^{1-a}F\in L^2(r\,dr\,d\theta)$, equivalently
$r^{2-a}F\in L^2(dr/r\,d\theta)$, then on a smaller disk
$u=u_0+\mathfrak h$, where $\mathfrak h$ is harmonic with Taylor series
containing only the allowed degrees, and
\[
  \int\Big(\frac{|u_0|^2}{r^{2a+2}}+\frac{|\nabla u_0|^2}{r^{2a}}
           +\frac{|D^2u_0|^2}{r^{2a-2}}\Big)\,dx\,dy<\infty .
\]
Moreover $u_0$ has a continuous representative with $|u_0|\le Cr^a$
(manuscript \texttt{(eq:subcritical-elliptic-norm)}).
\end{lemma}

\begin{proof}
Choose a radial cutoff $\chi$ supported in the original disk and equal to
one on a smaller disk. With $t=-\log r$, extend
$G=r^{2-a}\chi F$ by zero to the full cylinder $\R\times S^1$.
It is in $L^2(dt\,d\theta)$ by the hypothesis. For
$\mathcal P_a=(\partial_t-a)^2+\partial_\theta^2$, Fourier transformation
in $t$ and angular mode $m$ give
\[
 |(i\xi-a)^2-m^2|^2
 =\bigl(\xi^2+(m-a)^2\bigr)\bigl(\xi^2+(m+a)^2\bigr)
 \ge c_a(1+\xi^2+m^2)^2
\]
on the allowed modes. The positive constant exists because those modes
are discrete and none has $|m|=a$; for large $|m|+|\xi|$ the ratio is
bounded below directly. Thus inversion gives $V\in H^2$ on the full
cylinder, with $\mathcal P_aV=G$, preserving the reflection parity.
The Fourier Cauchy--Schwarz estimate gives $\|V\|_\infty\le C\|V\|_{H^2}$,
since $\sum_m\int_\R(1+\xi^2+m^2)^{-2}d\xi<\infty$.

Set $u_0=r^aV(-\log r,\theta)$. Polar differentiation gives the displayed
weighted norms and $|u_0|\le Cr^a$, while
$\Delta u_0=\chi F$ off the origin. It remains to exclude a distribution
at the origin. Let $\eta_\epsilon$ vanish on $r<\epsilon$ and equal one on
$r>2\epsilon$. On the intervening annulus $A_\epsilon$,
\[
 \|[\Delta,\eta_\epsilon]u_0\|_2
 \le C\epsilon^{a-1}\left(
   \|\nabla u_0/r^a\|_{L^2(A_\epsilon)}
  +\|u_0/r^{a+1}\|_{L^2(A_\epsilon)}\right)\longrightarrow0.
\]
The same estimates on the first and second derivatives show
$\eta_\epsilon u_0\to u_0$ in physical $H^2$ on the smaller disk; here
$a>1$ is essential. Hence $\Delta u_0=F$ distributionally there.
The difference $u-u_0$ is weakly harmonic, so Weyl's lemma makes it smooth
and harmonic across the origin. Its parity restricts the Taylor degrees.
\end{proof}

\section{Saddle estimates}

\begin{proposition}[First weighted vorticity estimate]
\label{prop:saddle-first-gain}
\uses{def:eigen-problem}
Every eigenfunction satisfies $\omega/d\in L^2$ (manuscript
\texttt{(prop:saddle-gains)}, first assertion).
\end{proposition}

\begin{proof}
\uses{lem:transport-first-gain, lem:transport-pair-eigen,
 lem:symmetries-commute}
Lift $(\omega,\psi)$ to bounded orbit neighbourhoods of the closed cells
and apply Lemma~\ref{lem:transport-first-gain}. For the negative cell use a
separatrix reflection, which preserves the eigenvalue equation and the
saddle distance. These estimates cover all torus saddle neighbourhoods;
away from them the conclusion follows from $\omega\in L^2$.
The four-character decomposition needed for arbitrary parity is included
in the transport lemma.
\end{proof}

\begin{proposition}[Odd--odd saddle estimates]
\label{prop:saddle-dd}
\uses{def:eigen-problem, def:sectors, def:sobolev-local}
In the odd--odd sector, $\omega/d^2\in L^2$, and at each saddle there is a
unique homogeneous harmonic quadratic $q=a\,uv$ (in the coordinates of
Lemma~\ref{lem:saddle-normal-form}) such that
\[
  \int_{r<r_0}\Big(\frac{|\psi-q|^2}{r^8}+\frac{|\nabla(\psi-q)|^2}{r^6}
  +\frac{|D^2(\psi-q)|^2}{r^4}\Big)<\infty
\]
(manuscript \texttt{(eq:quadratic-saddle-jet)}).
\end{proposition}

\begin{proof}
\uses{prop:saddle-first-gain, lem:cylinder-inverse, lem:transport-second-gain,
 lem:transport-pair-eigen, lem:saddle-normal-form, lem:stream-sector}
Fix $0<\delta<1$ and apply Lemma~\ref{lem:cylinder-inverse} with
$a=2-\delta$. Its physical forcing hypothesis is
$r^{\delta-1}\omega\in L^2$, supplied by the first gain. Odd--odd reflection
parity permits only angular degrees $2,4,\ldots$, all strictly larger than
$a$. Every term of the harmonic remainder therefore has the same finite
weighted norm as $u_0$. Consequently
\[
 |\psi|\le Cr^{2-\delta},\qquad
 \nabla\psi/r^{2-\delta}\in L^2
\]
near each saddle. The pointwise bound uses the cylinder inverse's bound
on $u_0$ and the harmonic remainder's quadratic vanishing.
Apply Lemma~\ref{lem:transport-second-gain} to $(\omega,\psi)$ on the cells;
away from the saddles its weighted forcing hypotheses follow from ordinary
$H^2$ regularity. This gives $\omega/d^2\in L^2$.

Now apply the cylinder inverse with $a=3$. Its physical hypothesis is
$\omega/r^2\in L^2$, equivalently $r^{-1}\omega\in L^2(dr/r\,d\theta)$,
and $3$ avoids all permitted angular degrees. Subtract the degree-two term
$q$ of the harmonic remainder. The remaining harmonic terms have degree
at least four, so they have finite weighted norm for $a=3$. The permitted
quadratic is $q=a_Puv$. If two quadratics worked, their difference would
have finite $\int|q|^2r^{-8}dx\,dy$, impossible for a nonzero homogeneous
quadratic. This proves uniqueness.
\end{proof}

\begin{remark}[Digest note: the saddle bootstrap and its weights]
The first transport gain, the subcritical cylinder inverse, and the second
transport gain establish the saddle jet without any $H^1$ regularity of
$\omega$. The two cylinder steps use different physical weights:
$r^{\delta-1}\omega$ for $a=2-\delta$, and $r^{-2}\omega$ for $a=3$.
The expression $r^{-1}\omega$ in the manuscript's last step is a cylinder
$L^2$ input, not a physical-area $L^2$ input.
The subsequent $H^1$ and logarithmically weighted transport arguments in
Propositions~\ref{prop:dd-separatrix-h1} and
\ref{prop:dd-quotient-domain} are separate steps; the saddle bootstrap
does not supply them.
\end{remark}

\begin{corollary}[Mixed saddle jet]
\label{cor:mixed-saddle-jet}
\uses{def:eigen-problem, def:sectors, def:sobolev-local}
\notready
In either mixed sector, at each saddle there is a unique permitted linear
harmonic $\ell$ such that
$\int\big(|\psi-\ell|^2r^{-6}+|\nabla(\psi-\ell)|^2r^{-4}+|D^2(\psi-\ell)|^2r^{-2}\big)<\infty$.
\end{corollary}

\begin{proof}
\uses{prop:saddle-first-gain, lem:cylinder-inverse}
Lemma~\ref{lem:cylinder-inverse} with $a=2$; central odd parity restricts the
angular frequencies to odd integers, so the multiplier has no zero.
\end{proof}

\section{Odd--odd sector: regularity across the separatrix}

\begin{lemma}[Weighted forcing supplied by the quadratic saddle jet]
\label{lem:dd-weighted-forcing}
\uses{def:eigen-problem, def:sectors, def:cell-arcs, def:sobolev-local}
For an odd--odd eigenfunction,
\[
 \psi/d^2,\qquad \nabla\psi/d^{3/2}\in L^2
\]
near the saddles. The continuous representative of $\psi$ vanishes on
every separatrix edge.
\end{lemma}

\begin{proof}
\uses{prop:saddle-dd, lem:stream-sector, lem:continuous-representative,
 lem:separatrix-lines}
Write $\psi=q+R$ near a saddle. The homogeneous quadratic satisfies
$q/r^2=O(1)$ and $\nabla q/r^{3/2}=O(r^{-1/2})$, both square integrable
in physical dimension two. The remainder estimates
$R/r^4,\nabla R/r^3\in L^2$ imply the required weaker weights by
multiplication by bounded powers of $r$.
Finally the stream function has odd parity under both reflections.
Its continuous representative has the same parity everywhere, since the
continuous parity identities hold almost everywhere. It therefore vanishes
on the fixed lines, which are exactly the separatrix.
\end{proof}

\begin{proposition}[$H^1$ gain near the separatrix, odd--odd]
\label{prop:dd-separatrix-h1}
\uses{def:eigen-problem, def:sectors, def:sobolev-local}
In the odd--odd sector, $\omega\in H^1$ on a neighbourhood of the separatrix
(including the saddles), with zero trace on each regular edge, and $\psi\in H^3$
on a possibly smaller neighbourhood (manuscript \texttt{(prop:separatrix-h1)}).
\end{proposition}

\begin{proof}
\uses{lem:dd-weighted-forcing, lem:weighted-transport-h1,
 lem:transport-zero-edge-trace, lem:transport-pair-eigen,
 lem:h1-crossing-gluing, lem:elliptic-local, lem:symmetries-commute}
Apply Lemma~\ref{lem:weighted-transport-h1} to $(\omega,\psi)$.
Lemma~\ref{lem:dd-weighted-forcing} supplies its weighted hypotheses; away
from the saddles these are ordinary $H^2$ bounds. Thus $\omega$ is $H^1$
on the cellwise separatrix strips, with $\omega/d^2\in L^2$.
The same argument on adjacent cells follows by reflection.

The forcing $\psi$ is zero on every regular edge. By
Lemma~\ref{lem:transport-zero-edge-trace}, both one-sided traces of
$\omega$ are zero. ACL trace gluing therefore gives $H^1$ across the
regular edges. At a saddle, $\omega/r\in L^2$ follows from the stronger
$\omega/d^2$ bound, and Lemma~\ref{lem:h1-crossing-gluing} glues the four
quadrants including their common endpoint. Finitely many such patches
give an open neighbourhood of the whole separatrix.
Finally $\Delta\psi=\omega\in H^1$ there and $\psi\in H^2$ globally;
Lemma~\ref{lem:elliptic-local} gives $\psi\in H^3$ on a smaller neighbourhood.
\end{proof}

\begin{remark}[Digest note: the $H^1$ gain without a midpoint condition]
The quadratic saddle jet supplies the weaker weighted-forcing hypotheses
of Lemma~\ref{lem:weighted-transport-h1}, so the former generic $H^1$
transport sketch is unnecessary. The endpoint ray terms are controlled by
Lemma~\ref{lem:weighted-corner-trace}; numerator differentiation uses only
$H^2$ regularity of $\psi$, not a prior $H^1$ bound for $\omega$.
Zero edge trace then follows from the boundary transport equation and
weighted corner integrability, with no separate assumption about the
midpoint value $A(0)$. The saddle extension is the explicit cutoff argument
of Lemma~\ref{lem:h1-crossing-gluing}, rather than an unspecified
removability assertion.
\end{remark}

\begin{lemma}[Logarithmically weighted forcing, odd--odd]
\label{lem:dd-log-forcing}
\uses{def:eigen-problem, def:sectors, def:sobolev-local, def:log-level-weight,
 def:cell-arcs, def:period}
For an odd--odd eigenfunction there is a neighbourhood $V$ of the
separatrix such that
\[
 \int_V\Lambda_\Psi\left(
       \frac{|\nabla\psi|^2}{d^3}+|D^2\psi|^2+|\nabla\psi|^2\right)<\infty.
\]
On each endpoint ray of a quarter orbit, near its saddle, the traces satisfy
\[
 \int_0^{r_0}(1+|\log r|)\left(
          \frac{|\psi|^2}{r^3}+\frac{|\nabla\psi|^2}{r}\right)dr<\infty.
\]
Gradient values on rays mean Sobolev traces, not arbitrary point values.
\end{lemma}

\begin{proof}
\uses{prop:dd-separatrix-h1, prop:saddle-dd, lem:saddle-normal-form,
 lem:log-multiplier-rectangle, lem:weighted-corner-trace,
 lem:h1-crossing-gluing}
Shrink the separatrix neighbourhood so that $\psi\in H^3$ there.
Near a saddle, the exact factorization $\Psi=uv\,m(u,v)$ gives
\[
 \Lambda_\Psi\le C\bigl(1+|\log|u||+|\log|v||\bigr).
\]
On a regular edge rectangle there is just one logarithmic coordinate
singularity. Lemma~\ref{lem:log-multiplier-rectangle}, localized on finitely
many rectangles, therefore controls $\int\Lambda_\Psi|F|^2$ for local
$H^1$ functions $F$. Apply this first to $D^2\psi$ and $\nabla\psi$,
which are $H^1$ by the ordinary $H^3$ gain.

For the remaining term write $\psi=q+R$ near a saddle. The quadratic part
has $|\nabla q|^2/r^3\le C/r$. Its logarithmically weighted integral is
finite: in polar coordinates the radial integrals of $1+|\log r|$ and
the angular integrals of $|\log\sin\theta|+|\log\cos\theta|$ are finite.
For $F=\nabla R/r^{3/2}$ the saddle jet gives
\[
 F/r\in L^2,\qquad
 |\nabla F|\le C\left(|D^2R|/r^{3/2}+|\nabla R|/r^{5/2}\right)\in L^2.
\]
It is $H^1$ off the origin; its traces match across the axes there.
The cutoff argument of Lemma~\ref{lem:h1-crossing-gluing} extends it in
$H^1$ across the origin. Apply the logarithmic multiplier to each localized
component of $F$. Away from saddle neighbourhoods, $d$ is bounded below,
so ordinary $H^3$ and the same multiplier suffice.

For the ray estimates, apply Lemma~\ref{lem:weighted-corner-trace} to
$R/r^3$ and each component of $\nabla R/r^2$. Their quotients by $r$
and their gradients are $L^2$ by the quadratic jet, giving
\[
 \int_{\rm ray}|R|^2r^{-7}dr<\infty,\qquad
 \int_{\rm ray}|\nabla R|^2r^{-5}dr<\infty.
\]
These imply the desired weaker ray weights because
$r^4(1+|\log r|)$ is bounded. The quadratic contributions are bounded by
$Cr(1+|\log r|)$ and are integrable directly.
\end{proof}

\begin{lemma}[Logarithmic numerator bounds, odd--odd]
\label{lem:dd-log-numerators}
\uses{def:eigen-problem, def:sectors, def:numerator, def:log-level-weight,
 def:cell-arcs, def:sobolev-local}
For an odd--odd eigenfunction with fixed rotation character
$\kappa\in\{1,-1\}$, let $N$ be its quarter numerator and put
\[
 B(h)=\int_{-L(h)}^0\psi(\gamma_h(t))e^{-i\lambda t}dt,\qquad
 K(h)=\int_{-L(h)}^{L(h)}\frac{|\nabla\psi(\gamma_h(t))|^2}{d^3}dt.
\]
For sufficiently small $h_0>0$,
\[
 \int_0^{h_0}\Lambda(h)\left(K(h)+|N'(h)|^2+|B'(h)|^2
                                  +h^2|N''(h)|^2\right)dh<\infty.
\]
\end{lemma}

\begin{proof}
\uses{lem:dd-log-forcing, lem:dd-weighted-forcing,
 lem:weighted-numerator-estimates, lem:transport-pair-eigen,
 lem:orbit-coordinates-area, lem:period-orbit}
Apply Lemma~\ref{lem:weighted-numerator-estimates} to $(\omega,\psi)$;
Lemma~\ref{lem:dd-weighted-forcing} supplies its hypotheses and hence its
weak differentiations. With $p_\pm=\gamma_h(\pm L(h))$, its pointwise
estimates are
\begin{align*}
 |N'|^2+|B'|^2&\le CK(h)+C\sum_\pm|\psi(p_\pm)|^2/h^2,\\
 h^2|N''|^2&\le C\int_{-L}^L(|D^2\psi_h|^2+|\nabla\psi_h|^2)dt
       +C\sum_\pm\left(|\nabla\psi(p_\pm)|^2/h
                                +|\psi(p_\pm)|^2/h^2\right).
\end{align*}
Multiply by $\Lambda(h)$. Weighted coarea controls all interior terms by
Lemma~\ref{lem:dd-log-forcing}. On each endpoint ray,
$h=\sin^2(r/2)$, $dh\asymp r\,dr$ and
$\Lambda(h)\asymp1+|\log r|$, so the ray estimates in that lemma control
both endpoint terms. This proves the asserted weighted integrability.
\end{proof}

\begin{proposition}[Weighted vorticity gradient, odd--odd]
\label{prop:dd-weighted-gradient}
\uses{def:eigen-problem, def:sectors, def:cell-arcs, def:sobolev-local}
For an odd--odd eigenfunction, $\nabla\omega/d\in L^2$ on a sufficiently
small neighbourhood of the separatrix.
\end{proposition}

\begin{proof}
\uses{lem:dd-log-numerators, lem:log-phase-cell-division, lem:numerator-zeros,
 prop:dd-separatrix-h1, prop:saddle-dd, lem:rotation-characters,
 lem:transport-unknown, lem:weighted-numerator-estimates,
 lem:orbit-level-derivatives, lem:orbit-coordinate-gradients,
 lem:orbit-coordinates-area, lem:period-coarse, lem:period-derivative-bounds}
Work with either rotation character and recenter time separately on each
quarter patch. The numerator vanishes at divisor zeros, and
Lemma~\ref{lem:log-phase-cell-division} gives
\[
 \int_0^{h_0}\Lambda(h)\bigl(|C'|^2+|C/h|^2\bigr)dh<\infty.
\]
In the midpoint identity
$A=ie^{2i\lambda L}C+iB$, use $|L'|\le C/h$ and
Lemma~\ref{lem:dd-log-numerators} to obtain
$\int\Lambda|A'|^2dh<\infty$.
For $z=-(\omega+\psi)/\lambda$, midpoint reconstruction gives
\[
 \partial_hz_h(t)=e^{i\lambda t}A'(h)
       +i\int_0^te^{i\lambda(t-u)}\partial_h\psi_h(u)du.
\]
The subinterval estimate of Lemma~\ref{lem:weighted-numerator-estimates}
bounds the squared integral by $CK(h)$. Since $2L(h)\le C\Lambda(h)$,
\[
 \int_0^{h_0}\int_{-L(h)}^{L(h)}|\partial_hz_h|^2dt\,dh
 \le C\int_0^{h_0}\Lambda(h)(|A'|^2+K(h))dh<\infty.
\]
The jet also gives $\nabla\psi/d\in L^2$, so
$|\partial_h\psi_h|\le C|\nabla\psi_h|/d$ proves
$\partial_h\omega_h\in L^2(dh\,dt)$.

For the time derivative,
\[
 \partial_t\omega_h=i\lambda\omega_h-\partial_t\psi_h,\qquad
 |\partial_t\psi_h|\le Cd|\nabla\psi_h|.
\]
Thus $d^{-2}\partial_t\omega_h\in L^2$, by
$\omega/d^2\in L^2$ and $\nabla\psi/d\in L^2$.
The inverse-coordinate estimates give
\[
 \frac{|\nabla\omega|}{d}
 \le C\left(|\partial_h\omega_h|+
                        d^{-2}|\partial_t\omega_h|\right).
\]
Coarea now proves the claim on each quarter strip. Sum the character
components and use reflection for the adjacent cells. The weak gradient
is the one already supplied across edges and saddles by
Proposition~\ref{prop:dd-separatrix-h1}; no new jump term is introduced.
\end{proof}

\begin{lemma}[$L^2$ vorticity quotient, odd--odd]
\label{lem:dd-vorticity-quotient-l2}
\uses{def:eigen-problem, def:sectors, def:sobolev-local}
For an odd--odd eigenfunction, $\nu=\omega/\Psi\in L^2(\T^2)$.
\end{lemma}

\begin{proof}
\uses{prop:dd-weighted-gradient, prop:dd-separatrix-h1,
 lem:saddle-normal-form, lem:hardy-1d, lem:acl-slicing, lem:trace-l2}
Near a saddle use a small quadrant disk $Q$ in the orthonormal coordinates.
The exact
factorization gives $|\Psi|\asymp r^2\sin\theta\cos\theta$ there.
The $H^1$ vorticity has zero traces on both axes, so polar ACL slicing on
annuli and the angular Hardy inequality give, for almost every $r$,
\[
 \int_0^{\pi/2}\frac{|\omega(r,\theta)|^2}
                         {\sin^2\theta\cos^2\theta}\,d\theta
 \le C\int_0^{\pi/2}|\partial_\theta\omega(r,\theta)|^2d\theta.
\]
Multiply by $r^{-3}dr$ and use
$|\partial_\theta\omega|\le r|\nabla\omega|$ to get
\[
 \int_Q|\omega/\Psi|^2\le C\int_Q|\nabla\omega|^2/r^2<\infty.
\]
On regular edge rectangles, $|\Psi|$ is comparable to distance to the edge;
zero vorticity trace and endpoint Hardy give the quotient in $L^2$.
Away from the separatrix, division is bounded. A finite cover gives the
claim on the torus.
\end{proof}

\begin{remark}[Digest note: the logarithmic gradient gain]
The logarithm pays for the growing quarter length when controlling
$\partial_h\omega$ without the factor $d$ present in the ordinary $H^1$
gain. It is not differentiated: phase-cell comparability transfers the
existing Hardy estimates to $\Lambda(h)dh$.
The manuscript uses $H^1\hookrightarrow L^4$ to control multiplication by
$\Lambda_\Psi$. The rectangle slicing argument of
Lemma~\ref{lem:log-multiplier-rectangle} supplies exactly that multiplier
bound without a finite-$p$ embedding. This resolves the weighted gradient
and $L^2$ vorticity division steps, but does not establish the $H^2$ stream
quotient or its weighted corner remainders.
\end{remark}

\begin{proposition}[Quotient domain, odd--odd]
\label{prop:dd-quotient-domain}
\uses{def:eigen-problem, def:sectors, def:sobolev-local}
\notready
For every odd--odd eigenfunction,
\[
  \phi:=\psi/\Psi\in H^2(\T^2),\qquad \nu:=\omega/\Psi\in L^2(\T^2),\qquad
  \nu/d\in L^2(\T^2).
\]
$\phi$ is even under both separatrix reflections, $(\nu,\phi)$ is a transport
pair ($H(\nu+\phi)=\lambda\nu$ in the graph domain), and at each saddle there is
a constant $a$ with $(\phi-a)/r^2,\ \nabla\phi/r\in L^2$ locally and
$\phi(\mathrm{saddle})=a$ (manuscript \texttt{(prop:dd-quotient-domain)}).
\end{proposition}

\begin{proof}
\uses{prop:dd-separatrix-h1, prop:saddle-dd, prop:dd-weighted-gradient,
 lem:dd-vorticity-quotient-l2, lem:transport-first-gain, lem:hardy-1d,
 lem:elliptic-local, lem:sobolev-holder, lem:local-to-global}
Proposition~\ref{prop:dd-weighted-gradient} and
Lemma~\ref{lem:dd-vorticity-quotient-l2} supply $\nabla\omega/d\in L^2$
near the saddles and $\nu\in L^2(\T^2)$.
The remaining proposed argument is as follows. Dyadic elliptic estimates give
$D^3R/r\in L^2$ for $R=\psi-a\Psi$, and the division formula
$R/(uv)=u^{-1}\int_0^1R_v(u,tv)\,dt$ gives the $H^2$ quotient. Multiplying
the graph equation by $\chi_\epsilon(\Psi)/\Psi$ (using $L_0\Psi=0$) and
closedness give the quotient transport equation. Finally
Lemma~\ref{lem:transport-first-gain} for $(\nu,\phi)$ gives $\nu/d\in L^2$.
\end{proof}

\section{Mixed sectors: regularity across the separatrix}

\begin{proposition}[Mixed-sector regularity]
\label{prop:mixed-quotient-domain}
\uses{def:eigen-problem, def:sectors, def:cells, lem:cell-square, def:sobolev-local}
\notready
In the mixed sector $(-,+)$ (odd across $a=\pm\pi/2$, even across
$b=\pm\pi/2$), $\omega\in H^1$ near the separatrix with zero trace on the odd
edges, $\psi\in H^3$ near the separatrix, and on the positive cell
\[
  \phi=\psi/\cos a\in H^2(\Cell),\qquad \nu=\omega/\cos a\in L^2(\Cell),
\]
including the corners (manuscript \texttt{(prop:mixed-quotient-domain)}).
\end{proposition}

\begin{proof}
\uses{cor:mixed-saddle-jet, lem:saddle-normal-form,
  lem:weighted-transport-h1, lem:elliptic-local, lem:hardy-1d,
  lem:weighted-corner-trace}
Choose a smooth $z_0$, affine near each saddle, whose linear term solves
$(H_s-\lambda)z_0=\ell$ (possible since $H_s$ has eigenvalues $\pm i$ on linear
functions). Then $(\omega_0,\psi_0)=(-Hz_0,(H-\lambda)z_0)$ is a transport pair.
The cubic Taylor remainder of the dipole field gives
$\psi_0=\ell+O(r^3)$ with the corresponding derivative bounds.
The remainder pair $(v,f)=(\omega-\omega_0,\psi-\psi_0)$ therefore has zero jet:
$f/d^3,\nabla f/d^2,D^2f/d\in L^2$ near the saddles.
In particular $f/d^2,\nabla f/d^2\in L^2$.
Lemma~\ref{lem:weighted-transport-h1} gives $v/d^2\in L^2$ and
cellwise $H^1$, hence cellwise $\omega\in H^1$.
On an odd edge the boundary trace satisfies
$-i\partial_t\omega=\lambda\omega$, hence equals $ce^{i\lambda t}$.
Since $\omega_0=O(d)$ and $v/d^2\in L^2$, one has $\omega/d\in L^2$;
Lemma~\ref{lem:weighted-corner-trace} forces $c=0$.
The odd traces are zero and the even traces agree by reflection, so
gluing and removability of the saddle points give $H^1$ across the separatrix.
Local elliptic regularity gives $\psi\in H^3$ near the separatrix.
Normal Hardy division of the zero-trace $H^1$ vorticity gives $\nu\in L^2$.
For the odd $H^3$ stream function, differentiate
$\psi(u,v)/u=\int_0^1\psi_u(tu,v)\,dt$ through order two;
Minkowski uses only the integrable dilation factor $t^{-1/2}$.
The smooth factor between $u$ and $\cos a$ gives the $H^2$ quotient.
Away from the odd edges, the original $\psi\in H^2$ suffices,
including at the elliptic center.
\end{proof}

\begin{remark}[The mixed quotient is not a transport pair]
Unlike $\Psi$, the divisor $h=\cos a$ is not invariant under the flow:
$L_0h=uvh$, where $u=\sin a$ and $v=\sin b$.
Consequently the mixed quotients satisfy
\[
 H(\nu+\phi)=\lambda\nu+iuv(\nu+\phi),
\]
not the transport-pair equation. The gains above are applied to the
jet-subtracted pair before division, not directly to $(\nu,\phi)$.
\end{remark}

\section{Elliptic centers}

\begin{proposition}[Regularity at the elliptic centers]
\label{prop:center-regularity}
\uses{def:eigen-problem, def:sobolev-local}
\notready
Every eigenfunction has $\omega\in H^1$ near each elliptic center, including
the case $\lambda\in\Z$ (manuscript \texttt{(prop:center-regularity)}).
\end{proposition}

\begin{proof}
\uses{lem:period-regularity, lem:period-elliptic, lem:elliptic-local, lem:hardy-1d}
Construct smooth action-angle-type coordinates $F(p)=\frac1{2\pi}\int_0^{2\pi}R_{-t}\Phi^t(p)\,dt$
for the normalized flow $L_0/b(\Psi)$, in which $H=b(r)D_\theta$ and
$b(r)=1-r^2/8+O(r^4)$. In angular modes
$\omega_n=-\frac{nb}{nb-\lambda}\psi_n$. For $\lambda\notin\Z$ the multipliers
are bounded with radial derivatives $O(r)$. For $\lambda=m\in\Z\setminus\{0\}$
the resonant mode solves $(\Delta_0-8/r^2)p=G$ with $G/r\in L^2$; solve on the
cylinder (indicial roots $\pm\sqrt{m^2+8}$) to get $p/r^3,\nabla p/r^2,D^2p/r\in L^2$.
\end{proof}

\begin{corollary}[Global $H^1$ regularity]
\label{cor:global-h1}
\uses{def:eigen-problem, def:sectors, def:sobolev-local}
\notready
In the odd--odd and mixed sectors every eigenfunction has $\omega\in H^1(\T^2)$
and $\psi\in H^3(\T^2)$ (manuscript \texttt{(cor:global-eigenfunction-h1)}).
\end{corollary}

\begin{proof}
\uses{prop:dd-separatrix-h1, prop:mixed-quotient-domain, prop:center-regularity, lem:resonant-vanishing, lem:hardy-1d, lem:elliptic-local}
Separatrix neighbourhoods and centers are covered above. On compact regular
annuli only finitely many angular modes can be resonant. For each, the
numerator is $H^2$ in the transverse variable and vanishes at the simple
resonant level, and division gives an $H^1$ quotient. The other modes have
uniformly bounded multipliers.
\end{proof}

\begin{remark}[Digest note: elliptic centers are not a sector prerequisite]
The manuscript's only use of Corollary~\ref{cor:global-h1} in a sector
exclusion is the mixed corner trace. Lemma~\ref{lem:weighted-corner-trace}
supplies that trace using local $H^1$ and $\omega/r\in L^2$, both already
in Proposition~\ref{prop:mixed-quotient-domain}; see
Lemma~\ref{lem:mixed-boundary-square-traces}.
Thus the global corollary and Proposition~\ref{prop:center-regularity},
including its resonant case $\lambda\in\Z$, are separate assertions rather
than prerequisites of either remaining exclusion. Their sketches here
are retained as unverified arguments, not used to justify the local route.
\end{remark}
